\documentclass[tikz,border=10pt]{standalone}
\usetikzlibrary{shapes.geometric, arrows, positioning}

\usepackage[utf8]{inputenc} % Modifier selon l'encodage (latin1, utf8...)
\usepackage[T1]{fontenc}
\usepackage[french]{babel}
\usepackage{calligra}
\usepackage{bera}

\usepackage{amsmath}          % Permet de taper des formules mathématiques
\usepackage{amssymb}          % Permet d'utiliser des symboles mathématiques
\usepackage{amsfonts}         % Permet d'utiliser des polices mathématiques
\usepackage{amsthm}           % Permet d'écrire des théorèmes
\usepackage{nicefrac}         % Fractions 'inline'
\usepackage{stmaryrd}         % Pour les brackets
\usepackage{mathtools}        % Pour la KL div
\usepackage{bm}               % Pour les vecteurs en gras
\usepackage{bbm}

\begin{document}

\begin{tikzpicture}[
    node distance=2cm and 3cm,
    every node/.style={font=\small},
    block/.style={rectangle, draw, minimum width=2.5cm, minimum height=1.2cm, align=center},
    arrow/.style={draw, thick, ->, >=stealth}
]

% Nodes
\node [fill=yellow!30] (agent) at (0,2) [block] {Agent};
\node [fill=yellow!30] (env) at (0,-2) [block] {Environment};
\node (state_t) at (3,0.5) {State $S_t$};
\node (reward_t) at (3,-0.5) {Reward $R_t$};
\node (state_t1) at (-3,0.5) {State $S_{t+1}$};
\node (reward_t1) at (-3,-0.5) {Reward $R_{t+1}$};

% Arrows
\draw[arrow] (state_t) |- (agent);
\draw[arrow] (agent) -| (state_t1);
\draw[arrow] (reward_t1) |- (env);
\draw[arrow] (env) -| (reward_t);

\end{tikzpicture}

\end{document}
